% Provenance: copied from the design workspace CambridgePQ/recursion.fast, `report/sections/12-baseline.tex`, on 2026-10-19. Read-only copy; the source is authoritative.
% 12-baseline -- rewritten by W9 from submissions/baseline-v03/RESULT.md (P3b, contract v0.3 + round 3),
% submissions/baseline/RESULT.md (P3, contract v0.2) and the orchestrator's re-checks
\section{Baseline and MVP}\label{sec:baseline}
\begingroup\emergencystretch=3em

This section reports the first complete submission: the trivial ``proof $=w$'' verifier, compiled to a
RISC-V image and certified in Lean. It is certified twice. The \emph{current} certificate is against
the final contract v0.3 with the round-3 fixes (\cref{sec:cost}; contract \repo{RESULT.md} §R3), and it
passes the production harness. The \emph{historical} first certificate, against contract v0.2, is kept
intact. The section answers three questions. Is the certificate of
\cref{sec:protocol,sec:security,sec:cost} satisfiable by a real image? What does a complete
submission cost to write and to check? And what did writing one reveal about the contract?

Every number below is traceable to a file:
\begingroup\raggedright
\begin{itemize}[nosep]
\item {}[P3b~§n] is the v0.3 report \repo{submissions/baseline-v03/RESULT.md}. Its logs are in
  \repo{submissions/baseline-v03/logs/}: \repo{harness-T1.log}, \repo{harness-L1.log},
  \repo{check.log}, \repo{scores.log}, \repo{leanchecker.log}, \repo{loc.log} and \repo{emu.log}.
  Files named \code{L1} (\repo{harness-L1.log}, \repo{claim-l1.json}, the stub \code{L1.lean.in})
  belong to the latency record, whose certificate is \lean{Latency1}; they have nothing to do with
  security level L1.
\item {}[P3~§n] is the v0.2 report \repo{submissions/baseline/RESULT.md}, with its logs.
\item {}[O-recheck] is an independent re-check of the v0.2 submission, 2026-10-11, by the study's
  coordinator. {}[O-recheck-v03] is the same re-run of the v0.3 check, logged in
  \repo{submissions/baseline-v03/logs/orch-recheck.log}. After the round-4 harness changes (gate
  computability scope, build binding) it was repeated once more: CHECK OK in 13\,min 31\,s, T1 and
  \lean{Latency1} both ACCEPT in 364\,s and 371\,s (\repo{orch-recheck-r4.log}).
\end{itemize}
\endgroup

\begin{takeaway}[Headline: the final T1 certificate and the latency record are inhabited, pass the production harness, and the bound is tight]
\begin{itemize}[nosep,leftmargin=*]
\item \emph{Certificates.} \lean{RecursionFastCandidate.certificate : Certificate1 submission declared}
  is proved with no \lean{sorry} for $k\in\{16,20,24\}$, every clause included (summary, score,
  security, admission, image size, \lean{Agrees}, the accepting-run bound on all 18 counters). The
  optional latency record \lean{Latency1} is also proved, with work $=$ span $\le19+548\cdot2^k$
  [P3b~§0].
\item \emph{Production harness.} Comparator first in a fresh tree, Lean kernel \emph{and} nanoda
  \cite{nanoda} accepting, then the gate with \code{--manifest} in a second fresh tree: \textbf{HARNESS:
  ACCEPT} for T1 and for the \lean{Latency1} record; every audited theorem uses only
  \lean{[propext, Classical.choice, Quot.sound]}. An independent re-run of \code{check.sh} gave
  \code{CHECK OK} in 13\,min 35\,s, the two harness runs taking 365.8\,s and 373.2\,s
  [P3b~§5; O-recheck-v03].
\item \emph{Tight bound.} The 18 claimed counters are exact closed forms in rows and nonzeros, and
  emulator runs at the nonzero cap equal them at $k=16,20,24$. The latency bound is looser:
  $548\cdot2^k$ proved against $376\cdot2^k$ measured.
\item \emph{Scores.} $\mathrm{Score}_{zk}\approx1.624\times10^9$ and
  $\mathrm{Score}_{\mathrm{lat}}\approx6.105\times10^8$; padding of the index oracle is about 39\% of the
  RISC~Zero score (I-BAS-2).
\item \emph{Cost.} 1{,}920 code lines ($+439$ over v0.2, about 170 of them for the latency clause). One
  harness run takes about 6\,min and the full warm check 13.6\,min, against the brief's 1\,h. One
  \lean{simp [m3]} drove nanoda out of memory at 60\,GB; \lean{exact m3} fixed it.
\end{itemize}
\end{takeaway}

%-----------------------------------------------------------------------------------------------
\subsection{What the baseline is, and why it is the right MVP}\label{subsec:base-what}

\paragraph{A non-vacuity witness for the whole certificate.} Before P3, the T1 certificate was
inhabited only in parts. The contract's controls prove the security clause for a toy ``send the
witness'' protocol at a test size (\lean{toy\_secure}), and a smoke run evaluates the machine at full
fuel. No image satisfied \lean{Agrees} together with the accepting-run cost clause at the class shapes
[contract RESULT §8, item 3]. A certificate that no real image can satisfy would be a vacuity bug of
the kind listed in \cref{subsec:prior-lessons}. The baseline closes the gap: the rendered challenge
is solvable, and the full pipeline runs end to end. It goes from the assembly source through the Lean
certificate, challenge rendering and the production harness to the scored claim. For v0.3 it is also
the first inhabitant of the latency clause and the first submission through the round-3 harness. It
is the reference point every later entry is measured against, and the starter submission agents copy
from.

\paragraph{The protocol.} At class exponent $k$ the class sizes are $m=2^k$ rows, $n_x=2^{10}$,
$n_w=2^k-1-2^{10}$ and $\mathsf{nnz}\le2^{k+4}$. The protocol is the trivial NP verifier:
\begin{itemize}[nosep]
\item \emph{One prover round, no challenges.} The single round is \lean{msg ⟨k, k⟩}: a raw vector of
  $2^k$ elements of $K$ in one leaf. Its first $n_w=2^k-1025$ entries are read as the witness, and the
  remaining 1{,}025 are ignored. The honest prover sends $w$, zero padded. The extractor reads $w$
  back, and $\eps=0$.
\item \emph{Holographic only in the degenerate sense.} \lean{VK i} is empty. The verifier sees the
  index only through the oracle \lean{IDX i}, which it \texttt{QUERY}s once, as a single leaf, and
  reads in full. This satisfies the contract's holographic interface
  (\cref{def:bg-holographic}): the verifier never receives $i$ in the clear. Its query complexity,
  however, is linear in the index.
\item \emph{A CSR-like index oracle that stores $z$-addresses.} \lean{IDX i} has shape
  $\langle k{+}6,k{+}6\rangle$, that is $2^{k+6}$ words in one leaf. For each row $r$ and each matrix
  $A,B,C$ it stores an \emph{end word} $\mathtt{IDXB}+8\cdot(\text{position after the segment})$,
  followed by the pairs $(\mathrm{addr}(z_c),M_{rc})$ of the segment's entries
  (\code{Encoding.lean}: \lean{seg}, \lean{rowsOf}, \lean{encRows}). The address is where the machine
  finds $z_c$ without copying anything:
  \begin{itemize}[nosep]
  \item address $0$, the image data word $1$, for the constant wire $c=0$;
  \item $\mathtt{TAPE\_BASE}+32+8(c-1)$ for the public input $x$ on the tape;
  \item $\mathtt{WB}+8(c-1-n_x)$ for the witness, which the message \texttt{QUERY} copies there.
  \end{itemize}
  The encoding uses $3m+2(|A|+|B|+|C|)\le3\cdot2^k+2\,\mathsf{nnz}\le35\cdot2^k$ of the
  $2^{k+6}=64\cdot2^k$ words.
\end{itemize}
Choosing addresses over column numbers is the index-preprocessing freedom of a holographic protocol
used in its simplest form. The indexer does the address arithmetic once, so the verifier performs
none per nonzero.

%-----------------------------------------------------------------------------------------------
\subsection{The spec protocol, the sparse indexer and security}\label{subsec:base-spec}

\paragraph{Segments over the sparse index.} In v0.3 an \lean{Index} is sparse: each matrix is an entry
list (\lean{SpMat}) with duplicates summed and the cap \lean{Index.nnz\_le}. This resolves the
baseline's v0.2 finding I-BAS-1 (\cref{subsec:base-findings}). The segment \lean{seg M r} is the sublist of
row-$r$ entries in list order. Duplicates and explicit zeros are kept, and the machine sums them
exactly as \lean{SpMat.mulVec} does. The indexer computes every row in one pass, a \code{foldr} with
\lean{Array.modify} in $O(m+\mathsf{nnz})$ (\lean{buckets}). Its correctness lemma is
\lean{getD\_buckets}. The CSR lemma no longer needs a dense view: \lean{dot\_seg} (a segment's sum is
\lean{M.mulVec z r}) is \lean{SpMat.foldl\_eq} plus \lean{List.map\_map} [P3b~§2]. The statements,
copied from \code{Encoding.lean} and \code{Protocol.lean} (proofs omitted):

\begin{leancode}
theorem getD_buckets (M : SpMat sz) (r : Fin sz.m) : (buckets M).getD r [] = seg M r

/-- **On honest oracles the spec verifier decides `R1CS`** (the CSR lemma). -/
theorem verify_honest (hk : k ≤ 26) (i : Index (csz k)) (x : PubInput (csz k))
    (tr : FullTranscript (shape k)) :
    verify k x (IDX k i) tr = decide (R1CS i x (witOf (nw_le k) (tr 0)))

/-- **Security of the baseline** at every `k ≤ 26`: perfect completeness, RBR knowledge
soundness with `ε = 0` (no challenges), and the error cap. -/
theorem baseline_secure (hk : k ≤ 26) : (baseline k).Secure
\end{leancode}

The spec verifier \lean{verify} is \lean{walkRows}. It follows the end words of the oracle exactly as
the machine does, which makes the refinement proof structural. \lean{verify\_honest} combines
\lean{walkRows\_encRows} (the walk over an honest encoding visits each row's three segments),
\lean{dot\_seg} with \lean{rowsOK\_rowsOf} (``all rows hold'' is \lean{R1CS}) and \lean{zAt\_zAddr}
(reading through the stored addresses yields $z=(1,x,w)$). \lean{rowsLen\_rowsOf} bounds the encoding
length by $3m+2\,\mathsf{nnz}$ from \lean{Index.nnz\_le}. The indexer is computable and linear by
construction, but it has not been run at class size; \lean{encRows} is a non-tail recursion over
$2^k$ rows [P3b~§2, §9].

\paragraph{Why RBR knowledge soundness is immediate.} The schedule has no challenge round, so
\lean{ChallengeIdx (shape k)} is empty (\lean{noChal}). The RBR-KS counting condition and the error
cap, both quantified over challenge indices, hold vacuously. The content therefore sits in the
knowledge state function \lean{ksf}, which the contract requires in any case:
\begin{itemize}[nosep]
\item before the message, the state is ``$R1CS(i,x,w)$'';
\item after the message, it is ``$R1CS(i,x,\mathsf{witOf}(\mathit{msg}))$'';
\item the \lean{full} clause says that an accepting transcript makes the second state true, and it
  follows from \lean{verify\_honest}.
\end{itemize}
The construction is the one in the contract's toy control. Vacuity is not a risk here: the contract
controls show that accept-all has no KSF on a one-message schedule (\lean{acceptAll\_no\_KSF},
\cref{tab:sec-controls}). Perfect completeness is \lean{verify\_honest} applied to the honest
transcript, with \lean{witOf\_msgOf}. Apart from the sparse segments, the security proof is unchanged
from v0.2.

%-----------------------------------------------------------------------------------------------
\subsection{The RISC-V image}\label{subsec:base-image}

The image is unchanged since v0.2 and byte-identical [P3b~§0]. The source is \repo{asm/baseline.S}.
\repo{asm/build.sh} assembles it with \code{riscv64-unknown-elf-as -march=rv64im\_zbc} into
\repo{baseline.bin} and \repo{baseline.data}, and generates \code{Image.lean} and \code{Code.lean} from
those raw files. For v0.3, \repo{asm/gen\_image.py} emits \code{Image.lean} with the contract
renderer's own chunked form, so that \lean{Challenge.image = image} holds by \lean{rfl}.

The image has 84 instructions (336\,B) and one data word, the constant $1$ at address 0. It uses nine
instruction forms: \code{ld}, \code{addi}, \code{slli}, \code{xor}, \code{clmul}, \code{clmulh},
\code{beq}, \code{bne} and \code{ecall} \cite{RISCV-Bitmanip}. There is no \code{jal} and no store:
the two \texttt{QUERY}s are the only memory writes. The layout is as follows.
\begin{itemize}[nosep]
\item \emph{Prologue}, instructions 0--13. It issues two \texttt{QUERY}s, for \lean{IDX} to
  $\mathtt{IDXB}=2^{34}$ and for the message to $\mathtt{WB}=2^{33}$. It sets
  $\mathtt{s0}=\mathtt{0x1B}$ (the reduction constant, $x^{64}\equiv x^4+x^3+x+1$) and
  $\mathtt{s1}=m$ (tape word 0).
\item \emph{Row loop.} Three copies of a 19-instruction segment, at instructions 14, 33 and 52,
  compute $\langle A_r,z\rangle$, $\langle B_r,z\rangle$ and $\langle C_r,z\rangle$. They are followed
  by the product, one reduction and \code{bne} to the label \code{reject}.
\item \emph{Rejection} branches to the address just past the last instruction. The fetch fails there,
  which is a rejection.
\item \emph{Acceptance} is \code{HALT(1)} at instruction 83.
\end{itemize}
The inner products are accumulated \emph{unreduced}, as the XOR of the 128-bit carry-less products
(\code{clmul}/\code{clmulh}), and reduced once per segment:
\[
  \mathsf{red}(\mathit{lo},\mathit{hi})=\mathit{lo}\oplus\mathsf{clmul}(\mathit{hi},\mathtt{0x1B})
  \oplus\mathsf{clmul}(\mathsf{clmulh}(\mathit{hi},\mathtt{0x1B}),\mathtt{0x1B}).
\]
In characteristic 2, addition has no carries, so the accumulator needs no overflow invariant at any
segment length. A3's prime-field calibration probe needed one. The arithmetic bridge is
\lean{red\_halves} ($\mathsf{red}$ of the two halves of $T$ is \lean{K.reduce T}) together with
\lean{red\_acc} (lazy accumulation reduces to $\sum v\cdot z$) in \code{Arith.lean}.

\begin{Verbatim}[fontsize=\small,frame=single,framerule=0.3pt,rulecolor=\color{black!30},
  framesep=4pt,label={\footnotesize asm/baseline.S: the segment macro (comments annotated)}]
.macro SEG res, L, D
    ld    t1, 0(s2)            # segment end (absolute end word from IDX)
    addi  s2, s2, 8
    addi  a3, zero, 0          # lo := 0
    addi  a4, zero, 0          # hi := 0
    beq   s2, t1, \D           # empty segment
\L:
    ld    t2, 0(s2)            # address of z[col], precomputed by the indexer
    ld    t3, 8(s2)            # value M[r][col]
    ld    t4, 0(t2)            # z[col]: data word, tape (x) or QUERY destination (w)
    clmul  t5, t3, t4          # low half of the 128-bit carry-less product
    clmulh t6, t3, t4          # high half
    xor   a3, a3, t5           # unreduced accumulation (no carries)
    xor   a4, a4, t6
    addi  s2, s2, 16
    bne   s2, t1, \L           # 9 instructions per nonzero
\D:
    clmulh a5, a4, s0          # reduce (lo, hi); s0 = 0x1B
    clmul  a5, a5, s0
    xor    a3, a3, a5
    clmul  a5, a4, s0
    xor    \res, a3, a5        # res := <M_r, z> in K
.endm
\end{Verbatim}

Control flow depends on the data in exactly one way: the early exit of a failing row. With $m$ rows
and $P$ nonzeros, an accepting run takes $14+40m+9P+3$ steps (\lean{rowsSteps\_eq}). At $k=24$, at the
cap, that is $3.09\times10^9$ steps, 4.5\% of the fuel $2^{36}$ [P3~§2; P3b~§3].

%-----------------------------------------------------------------------------------------------
\subsection{Machine-level theorems: Agrees, cost and latency in one pass}\label{subsec:base-machine}

\paragraph{One reachability predicate on the contract's machine.} The v0.3 migration notes offered
two routes. One re-points the v0.2 proofs at the frozen v0.2 machine \lean{V02.run} and transfers
them through \code{Library/V02.lean}. The other re-bases the proofs on v0.3's own \lean{execute}. P3b
took the second route and obtained decision agreement, the 18-counter bound and the latency bound
from \emph{one} set of block proofs [P3b~§2, §8 item 3]. The central notion in \code{Rv/L0.lean}
extends v0.2's \lean{Reach} with an accounting invariant and a timer potential:

\begin{leancode}
/-- From `s`, `n` steps reach `s'` on the root hart: the counters grow by `d`, the
footprint invariant is kept, the timer potential grows by `e`, whatever the remaining
fuel. -/
def Reach (s : MachineState) (n : ℕ) (d : CostVec) (e : ℕ) (s' : MachineState) :
    Prop :=
  ∀ f (x : Acct) (sb : SB) (Φ : ℕ), Inv img F x → x.work = sb.time → PB sb Φ →
    ∃ x' sb', Inv img F x' ∧ x'.cost = x.cost + d ∧ x'.work = sb'.time ∧
      PB sb' (Φ + e) ∧
      execute img os (n + f) {} s sb x = execute img os f {} s' sb' x'
\end{leancode}

\lean{Inv img F x} is the \emph{footprint invariant}: every touched word and page lies in the code or
in the data footprint $F$, and every dirty page lies in $F$. \lean{PB sb Φ} is the \emph{timer
potential}: every register is ready by cycle $\Phi$ and the next issue slot is at most $2\Phi$. Each
step lemma carries the extra obligations. A load's or \texttt{QUERY}'s bytes must lie in $F$, while
instruction fetch is in the code footprint by construction, derived from \lean{fetch} succeeding. An
issue raises $\Phi$ by $\max(\mathit{lat},1+\mathit{bubble})$ (\lean{PB.step}): by 1 for ALU
instructions, 3 for \code{ld} and Zbc, 6 for a branch (taken or not) and $1+2^{\mathit{logLeaf}}$ for
a \texttt{QUERY}. \lean{Agrees} follows from the same \lean{run\_accept}/\lean{run\_reject} as the
bounds. The \lean{V02} transfer and \lean{Seq.run\_eq} served an intermediate, timer-free variant,
which was built and certified first. They are no longer needed.

\paragraph{The statements.} Copied from \code{RecursionFastCandidate/Main.lean} (proofs omitted, long lines
re-wrapped):

\begin{leancode}
/-- The instruction and QUERY counters of an accepting run with `m` rows and `P`
nonzero entries, at class `k` (first touch is bounded separately: `card_words`,
`card_pages`). -/
def closedCost (k m P : ℕ) : CostVec where
  add := 9 + 10 * m + P
  logic := 8 * m + 2 * P
  shift := 4
  clmul := 14 * m + 2 * P
  ld := 1 + 3 * m + 3 * P
  br := 5 * m + P
  sys := 1
  qCount := 2
  qLevels := 0
  qBytes := 8 * 2 ^ (k + 6) + 8 * 2 ^ k

theorem totalCost_eq (R : List (Row (csz k))) :
    totalCost k R = closedCost k R.length (pairs R)

def wordsBound (k : ℕ) : ℕ := 1068 + 2 ^ k + 2 ^ (k + 6)
def pagesBound (k : ℕ) : ℕ :=
  11 + 2 * ((8 * 2 ^ k + 1023) / 1024 + (8 * 2 ^ (k + 6) + 1023) / 1024)

theorem card_words {x : Acct} (hx : Inv image (foot k) x) :
    x.words.card ≤ wordsBound k
theorem card_pages {x : Acct} (hx : Inv image (foot k) x) :
    x.pages.card + x.dirty.card ≤ pagesBound k

/-- **The claimed bound at class `k`**: instruction counters at `m = 2^k` rows and
`P = 2^(k+4)` nonzero entries, first touch by `wordsBound`, `pagesBound`. -/
def costBound (k : ℕ) : CostVec :=
  { closedCost k (2 ^ k) (2 ^ (k + 4)) with
    mWords := wordsBound k, mPages := pagesBound k }

/-- Timer potential of the rows: `99` per row, `24` per nonzero entry. -/
theorem rowsLat_eq {sz : Sizes} (R : List (Row sz)) :
    rowsLat R = 99 * R.length + 24 * pairs R

/-- **The latency bound at class `k`** (single hart: work = span): `19 + 548·2^k`
cycles. -/
def latBound (k : ℕ) : ℕ := 19 + 548 * 2 ^ k

theorem run_accept (hk : k ≤ 26) (i : Index (csz k)) (x : PubInput (csz k))
    (tr : FullTranscript (shape k)) (h : rowsOK (rowsOf i) (zOf x tr)) :
    ∃ s y sb, execute image (oraclesOf i tr) FUEL {} (initState image (tapeOf i x tr))
        (SB.start 0) {} = some (y, sb, s) ∧
      y.cost = totalCost k (rowsOf i) ∧ Inv image (foot k) y ∧ y.work = sb.time ∧
      sb.time ≤ ePro k + rowsLat (rowsOf i) + 3

theorem run_reject (hk : k ≤ 26) (i : Index (csz k)) (x : PubInput (csz k))
    (tr : FullTranscript (shape k)) (h : ¬ rowsOK (rowsOf i) (zOf x tr)) :
    execute image (oraclesOf i tr) FUEL {} (initState image (tapeOf i x tr))
      (SB.start 0) {} = none

/-- **Decision agreement** of the image with the baseline spec verifier
(every `k ≤ 26`). -/
theorem baseline_agrees (hk : k ≤ 26) : (baseline k).Agrees image

/-- **Accepting-run cost bound**: every accepting run has counters `≤ costBound k`
(all 18). -/
theorem baseline_cost (hk : k ≤ 26) :
    (baseline k).AcceptingCostAtMost image (costBound k)

/-- **Accepting-run latency bound** (latency board): pipeline work and span
`≤ latBound k` (one hart: work = span = the clock). -/
theorem baseline_latency (hk : k ≤ 26) :
    (baseline k).AcceptingLatencyAtMost image (latBound k) (latBound k)
\end{leancode}

The contract fixes \lean{Agrees}, \lean{AcceptingCostAtMost} and \lean{AcceptingLatencyAtMost}
(\cref{sec:protocol,subsec:cost-vector,subsec:cost-lat}). They quantify over every index of the class
shape (with honest \lean{VK}/\lean{IDX}), every $x$ and every raw message. \lean{baseline\_agrees} is a
case split on \lean{rowsOK} using \lean{run\_accept}, \lean{run\_reject} and \lean{verify\_honest}. In
\lean{baseline\_cost} and \lean{baseline\_latency} the rejecting case is vacuous. In the accepting
case the instruction counters are rewritten with \lean{totalCost\_eq} and $P\le2^{k+4}$ is
\lean{pairs\_le}, from \lean{Index.nnz\_le}. The first-touch counters come from the invariant through
\lean{card\_words} and \lean{card\_pages}. The time comes from the potential, $\mathtt{ePro}\,k+99m+24P+3$.
All theorems are generic in $k\le26$, and the bound on $k$ is used only for address arithmetic
($\mathtt{IDXB}+8\cdot2^{k+6}\le2^{36}$).

\paragraph{The certificate records.} \code{Certificate.lean} writes the claim, the shape summaries,
the claimed $\mathit{zk4}=\mathrm{Score}_{zk}^{48}$ (a rational with a 446-digit numerator) and the
latency claim as literals. All of them are generated from the contract's \emph{Lean} score functions
(\repo{scripts/Eval.lean} $\to$ \repo{gen-claim.py}), with no Python mirror of the scoring. They are
tied to the proved bounds by \lean{claim\_eq : claim c = costBound (classExp c)} and
\lean{lat\_eq} (each \lean{fin\_cases c <;> rfl}). The score clauses hold by \lean{decide +kernel}. From
\code{Certificate.lean} (proofs omitted, long lines re-wrapped):

\begin{leancode}
/-- **The T1 certificate of the baseline**, for any record with the declared
data. -/
theorem certificate_of (D : Declared1) (hI : D.image = image) (hS : D.summary = summary)
    (hC : D.claim = claim) (hZ : D.zk4 = zk4) : Certificate1 submission D

/-- **The T1 certificate of the baseline.** -/
theorem certificate : Certificate1 submission declared :=
  certificate_of declared rfl rfl rfl rfl

/-- **The latency-board record of the baseline**, for any record with the declared
image. -/
theorem latency_of (D : Declared1) (hI : D.image = image) (L : Fin 3 → ℕ × ℕ)
    (hL : L = lat) (q : ℚ) (hq : q = lat3) : Latency1 submission D L q
\end{leancode}

The rendered challenge declares its own constant \lean{declared} with the same value. The solution is
therefore \lean{certificate\_of declared rfl rfl rfl rfl}, plus \lean{latency\_of declared rfl lat rfl
\_ rfl} for the latency record, so the rendered image, summary, claim and $\mathit{zk4}$ are definitionally the
candidate's [P3b~§2].

\paragraph{Per-block counters.} The closed forms are sums of six block constants [P3b~§2]:
\begin{itemize}[nosep]
\item prologue $\{\mathrm{add}\,7,\mathrm{shift}\,4,\mathrm{ld}\,1,q^{\mathrm{count}}\,2\}$, timer
  potential $16+2^{k+6}+2^k$ (the two \texttt{QUERY} bubbles);
\item per segment, the head $\{\mathrm{add}\,3,\mathrm{ld}\,1,\mathrm{br}\,1\}$ (potential 12) and the
  reduction $\{\mathrm{logic}\,2,\mathrm{clmul}\,3\}$ (potential 11);
\item per nonzero, the inner loop $\{\mathrm{add}\,1,\mathrm{logic}\,2,\mathrm{clmul}\,2,\mathrm{ld}\,3,\mathrm{br}\,1\}$
  (potential 24);
\item per row, the tail $\{\mathrm{add}\,1,\mathrm{logic}\,2,\mathrm{clmul}\,5,\mathrm{br}\,2\}$ (potential 30);
\item the halt $\{\mathrm{add}\,2,\mathrm{sys}\,1\}$ (potential 3).
\end{itemize}
A row has potential $3(12+11)+30=99$. At $m=2^k$ and $P=2^{k+4}$ the run's potential is
$19+2^{k+6}+2^k+99\cdot2^k+24\cdot2^{k+4}=19+548\cdot2^k$. As a consistency check, the coarse sum
$\mathrm{add}+\mathrm{logic}+\mathrm{shift}=13+18m+3P$ is v0.2's proved $\mathrm{alu}$ count. The
other v0.2 counters map one-to-one: v0.2's ctrl is v0.3's br, and load is ld. \Cref{tab:base-counters}
gives the claim per class.

\paragraph{First touch.} The data footprint of class $k$ has five byte regions: data word~0, tape
header word~0, $x$, the witness destination \texttt{WB} ($2^k$ words) and the index destination
\texttt{IDXB} ($2^{k+6}$ words). Code adds 42 doublewords (84 instructions) and one page. The pages are
the data page, the 9 tape pages and the two \texttt{QUERY} destinations, which are written and so
count twice, in and out. Hence $m_{\mathrm{words}}\le1068+2^k+2^{k+6}$ and
$m_{\mathrm{pages}}\le11+2(2^{k-7}+2^{k-1})$ for $k\ge7$. No layout change lowers this: a
\texttt{QUERY} touches its whole leaf, so the $2^{k+6}$-word index oracle, of which $35\cdot2^k$
words are used, costs $2^{k-1}$ pages in and $2^{k-1}$ out. Only a smaller index oracle would cut
it (I-BAS-2).

\begin{table}[t]
\centering\footnotesize
\setlength{\tabcolsep}{4pt}
\begin{tabular}{@{}llrrr@{}}
\toprule
Counter & Closed form ($m$ rows, $P$ nonzeros) & $k=16$ & $k=20$ & $k=24$ \\
\midrule
$n_{\mathrm{add}}$ & $9+10m+P$ & 1{,}703{,}945 & 27{,}262{,}985 & 436{,}207{,}625 \\
$n_{\mathrm{logic}}$ & $8m+2P$ & 2{,}621{,}440 & 41{,}943{,}040 & 671{,}088{,}640 \\
$n_{\mathrm{shift}}$ & $4$ & 4 & 4 & 4 \\
$n_{\mathrm{clmul}}$ & $14m+2P$ & 3{,}014{,}656 & 48{,}234{,}496 & 771{,}751{,}936 \\
$n_{\mathrm{ld}}$ & $1+3m+3P$ & 3{,}342{,}337 & 53{,}477{,}377 & 855{,}638{,}017 \\
$n_{\mathrm{br}}$ & $5m+P$ & 1{,}376{,}256 & 22{,}020{,}096 & 352{,}321{,}536 \\
$n_{\mathrm{sys}}$ & $1$ & 1 & 1 & 1 \\
$m_{\mathrm{words}}$ & $\le1068+2^k+2^{k+6}$ & 4{,}260{,}908 & 68{,}158{,}508 & 1{,}090{,}520{,}108 \\
$m_{\mathrm{pages}}$ & $\le11+2(\lceil2^{k+3}/2^{10}\rceil+\lceil2^{k+9}/2^{10}\rceil)$ & 66{,}571 & 1{,}064{,}971 & 17{,}039{,}371 \\
$q^{\mathrm{count}}$ & $2$ & 2 & 2 & 2 \\
$q^{\mathrm{bytes}}$ & $8\cdot2^{k+6}+8\cdot2^k$ & 34{,}078{,}720 & 545{,}259{,}520 & 8{,}724{,}152{,}320 \\
\midrule
$W=S$ (proved) & \lean{latBound k} $=19+548\cdot2^k$ & 35{,}913{,}747 & 574{,}619{,}667 & 9{,}193{,}914{,}387 \\
\bottomrule
\end{tabular}
\caption{The claimed bound at $m=2^k$, $P=2^{k+4}$, i.e.\ \lean{claim c = costBound (classExp c)} and
\lean{lat} [P3b~§0, §1; \code{Certificate.lean}]. The other seven counters ($n_{\mathrm{mul}}$,
$n_{\mathrm{div}}$, $n_{\mathrm{ldsub}}$, $n_{\mathrm{st}}$, $n_{\mathrm{stsub}}$, $n_{\mathrm{jmp}}$,
$q^{\mathrm{lev}}$) are 0. Every value is attained by the emulator at the nonzero cap
(\cref{tab:base-emu}).}\label{tab:base-counters}
\end{table}

\paragraph{Why the bound is tight.} The instruction counters of an accepting run depend only on $m$
and $P$, so every satisfiable index with exactly $2^{k+4}$ nonzeros attains them. The first-touch
bound is attained when every $x$ column is referenced. P3b extended the emulator \repo{emu/rfemu.c} to
v0.3: the 18 counters with the contract's \lean{kind}, first touch (a hash set of keys,
\texttt{QUERY} intervals and page bitmaps, code fetch included) and the dual-issue scoreboard. The
instances were generated at the class shape with the cap reached (6+5+5 nonzeros per row, random $K$
values, $C$ solved on the constant wire). \Cref{tab:base-emu} shows the result: at all three classes
the accepting counters equal \lean{costBound} \emph{field by field}, all 18 of them
[P3b~§3; \repo{logs/emu.log}]. The three rejecting variants all reject: one random row broken, one
witness entry of row~0 changed, and the last row broken. The corrupted witness rejects after 197
steps, and the broken last row 4 steps before the accepting run would halt.

A differential test (\repo{emu/Diff.lean}, \repo{emu/diff.py}) runs the contract's own \lean{runFull}
(Lean \lean{\#eval}) and \code{rfemu} on two tiny instances ($m=2$, 11 nonzeros). Instance~0 accepts
in both, with identical counters
\code{[40, 38, 4, 0, 0, 50, 40, 0, 0, 0, 21, 0, 1, 113, 7, 2, 0, 544]} and $W=S=340$. Instance~1
rejects in both [\repo{logs/diff.log}]. The seven nonzero instruction counters equal
\lean{closedCost} at $m=2$, $P=11$; for instance $\mathrm{add}=9+20+11=40$. The emulator itself is
unverified. It is differential-tested on these instances, timer and first touch included, and
cross-checked against the proved closed forms at all three classes.

\paragraph{The latency bound is sound but not tight.} The scoreboard measures
$W=S=376\cdot2^k+7$ against the proved $548\cdot2^k+19$, a ratio of 1.457. The potential method
charges the \code{ld}$\to$use and Zbc$\to$\code{xor} latencies serially, and every branch 6 cycles,
whereas the dual-issue scoreboard overlaps them. A tighter proof would thread exact scoreboard states
per block. The contract library's \lean{time\_blocks} with its slack of $L_{\max}=16$ per block would
be looser still here, because the blocks are short [P3b~§8 item 5].

\begin{table}[t]
\centering\footnotesize
\setlength{\tabcolsep}{4pt}
\begin{tabular}{@{}rrrrrrr@{}}
\toprule
$k$ & steps & $m_{\mathrm{words}}$ & $m_{\mathrm{pages}}$ & $W=S$ measured & $W=S$ proved & emu time \\
\midrule
16 & 12{,}058{,}641 & 4{,}260{,}908 & 66{,}571 & 24{,}641{,}543 & 35{,}913{,}747 & 1.2\,s \\
20 & 192{,}938{,}001 & 68{,}158{,}508 & 1{,}064{,}971 & 394{,}264{,}583 & 574{,}619{,}667 & 19.5\,s \\
24 & 3{,}087{,}007{,}761 & 1{,}090{,}520{,}108 & 17{,}039{,}371 & 6{,}308{,}233{,}223 & 9{,}193{,}914{,}387 & 347\,s \\
\bottomrule
\end{tabular}
\caption{Accepting emulator runs at the nonzero cap [P3b~§3; \repo{logs/emu.log},
\repo{logs/emu-timer.log}]. In every row all 18 counters equal the claim of \cref{tab:base-counters};
the instruction counters are the same as in v0.2. Measured $W=S$ is $376\cdot2^k+7$. Generating the
$k=24$ instance takes 263\,s; the emulator run needs 4.7\,GB.}\label{tab:base-emu}
\end{table}

%-----------------------------------------------------------------------------------------------
\subsection{Score}\label{subsec:base-score}

\paragraph{The recursion board.} The certified $\mathit{zk4}$ is the contract's \lean{scoreZk4} of the
declared summaries and claim, under the final pricing with the leaf hash charged as
$P_{H,P}\cdot(q^{\mathrm{bytes}}/64+q^{\mathrm{count}})$ (\cref{subsec:cost-zk}; contract §R3.2).
\Cref{tab:base-score} gives the twelve $S_{P,k}$. The schedule has one message round, no
challenges and blow-up $V_k=2^k/(2^k+2^{k+4})=1/17$, so the blow-up factor is 1. The profile and
v0.2-display values are also proved in Lean, as \lean{S\_<profile>\_<k>} and \lean{S\_v02\_<k>} in the
generated \code{Score.lean}, each by \lean{decide +kernel}.

\begin{table}[t]
\centering\footnotesize
\setlength{\tabcolsep}{4pt}
\begin{tabular}{@{}rrrrrr@{}}
\toprule
$k$ & ZK-A ZisK & ZK-B SP1-H & ZK-C RISC~Zero & ZK-D Jolt & v0.2 weights (display) \\
\midrule
16 & 52{,}780{,}165.25 & 111{,}204{,}032.21 & 147{,}427{,}411 & 122{,}841{,}941 & 29{,}624{,}405 \\
20 & 844{,}323{,}973.25 & 1{,}778{,}901{,}901.01 & 2{,}358{,}437{,}971 & 1{,}965{,}058{,}901 & 473{,}958{,}485 \\
24 & 13{,}509{,}024{,}901.25 & 28{,}462{,}067{,}801.81 & 37{,}734{,}606{,}931 & 31{,}440{,}530{,}261 & 7{,}583{,}303{,}765 \\
\midrule
per $2^k$ & $\approx805$ & $\approx1{,}696$ & $\approx2{,}249$ & $\approx1{,}874$ & $\approx452$ \\
geomean over $k$ & $8.44\times10^8$ & $1.78\times10^9$ & $2.36\times10^9$ & $1.97\times10^9$ & $4.74\times10^8$ \\
\midrule
\multicolumn{5}{@{}l}{$\mathrm{Score}_{zk}=\bigl(\prod_P\prod_kS_{P,k}\bigr)^{1/12}$ (certified as $\mathit{zk4}=\mathrm{Score}_{zk}^{48}$)} & $1.624\times10^9$ \\
\multicolumn{5}{@{}l}{$\mathrm{Score}_{\mathrm{lat}}=\bigl(\prod_k(W_k/16+S_k)\bigr)^{1/3}$, $W_k=S_k=$ \lean{latBound k}} & $6.105\times10^8$ \\
\bottomrule
\end{tabular}
\caption{Baseline scores under the final v0.3 pricing [P3b~§6; \repo{logs/scores.log}; \code{Score.lean}].
The last two rows of the upper block are derived here from the twelve values (per $2^k$ at $k=16$).
The v0.2 column evaluates the same counters with the v0.2 weights; it is a display, not part of
$\mathrm{Score}_{zk}$.}\label{tab:base-score}
\end{table}

\paragraph{Where the cost goes.} Per $2^k$ at $k=16$, RISC~Zero's total of about 2{,}250 splits
into [P3b~§6]:
\begin{itemize}[nosep]
\item first-touched pages, $1.016\times1{,}130\approx1{,}148$, almost all of it the index
  oracle read in and written out ($\approx1{,}130$);
\item \texttt{QUERY} leaf hashing, $P_H\cdot q^{\mathrm{bytes}}/64=100\cdot520/64\approx812$;
\item instructions, $\approx224$;
\item proof bytes, $\lceil B/8\rceil\approx65$.
\end{itemize}
The other profiles do not price pages, so their totals are lower. ZisK is the cheapest profile at
$\approx805$ per $2^k$ and RISC~Zero the dearest at $\approx2{,}249$.

\paragraph{The latency board.} With one hart, work equals span, and the board's
$T_k=W_k/16+S_k=\tfrac{17}{16}\,\mathtt{latBound}\,k$. That gives
$\mathrm{Score}_{\mathrm{lat}}=\bigl(\prod_kT_k\bigr)^{1/3}\approx6.105\times10^8$, certified as
\lean{lat3} by \lean{decide +kernel}. On the measured scoreboard time ($376\cdot2^k$ instead of
$548\cdot2^k$) the same formula would give about 1.46 times less. A tighter timing proof of the same
image can gain up to that factor.

\paragraph{Comparison with v0.2.} The v0.2 certificate's $\mathrm{Score}_{\mathrm{T1}}=4.74\times10^8$
(\cref{subsec:base-v02}) is not comparable with $\mathrm{Score}_{zk}$. On the same counters the v0.2
weights give $S_{16}=29{,}624{,}405$, $1{,}040$ above v0.2's certified value: $1{,}024$ for absorbing
$x$ (A4-F6) and $16=P_H\,q^{\mathrm{count}}$ for the leaf-hash rounding [P3b~§6]. The v0.3 profiles
cost 1.8--5.0 times the v0.2 display, mostly for first touch, hash prices and per-instruction weights.

\begin{issue}{I-BAS-2}{Power-of-two oracle lengths and whole-leaf charging penalise padding (register: M-8)}
The index needs at most $35\cdot2^k$ words, but oracle lengths are powers of two, so its shape is
$2^{k+6}$ words, and a \texttt{QUERY} charges its whole leaf. Under the v0.2 weights the unused
$29\cdot2^k$ words were about 13\% of $S$ [P3~§6]. Under the final v0.3 pricing they are paid as
first-touched pages in and out ($0.45\cdot2^k$ pages $\times1{,}130\approx512$ per $2^k$ under
RISC~Zero), as leaf hashing ($P_H\cdot29\cdot8/64\approx362$ per $2^k$), and once more in the
proof-byte charge. The first two alone, 874 of about 2{,}250 per $2^k$, make padding
\textbf{about 39\% of the RISC~Zero score} [P3b~§6, §8 item 2]. A submitter can recover most of
it with a fixed \texttt{QUERY} loop over 35 leaves of $2^k$ words (index shape $\langle k{+}6,k\rangle$),
which cuts the index pages and leaf bytes by $29/64$ for about $35\cdot6$ Merkle levels. P3b estimates
the prologue loop proof at 100--150 lines and did not do it, keeping the image unchanged. We recommend
the contract fix instead: allow \lean{OShape} lengths that are multiples of the leaf rather than
powers of two, so that the cost model does not reward layout engineering. Tracked as M-8 in the issue
register (\cref{subsec:iss-open}).
\end{issue}

%-----------------------------------------------------------------------------------------------
\subsection{Proof engineering}\label{subsec:base-proof}

\paragraph{Architecture.} Everything in the submission is untrusted and kernel-checked.
\begin{itemize}[nosep]
\item \code{Rv/L0.lean} (339 lines; 171 in v0.2) is independent of the program. It defines
  \lean{Reach} (\cref{subsec:base-machine}) with \lean{Reach.trans}, so that fuel, the 18 counters, the
  footprint invariant and the timer potential compose and no other file reasons about fuel.
\item \code{Rv/L0.lean} also provides one lemma per instruction form with \emph{symbolic} registers and
  immediates (\code{addi}, \code{slli}, \code{xor}, \code{ld} with its alignment and range side
  conditions, Zbc, \code{beq}, \code{bne}, and \texttt{QUERY} with the closed-form memory), together
  with lemmas for \code{HALT} and fetch failure. Each now also emits the v0.3 ticks, the footprint
  obligation and the potential step.
\item \code{Code.lean} is generated. It proves \lean{code\_j : CodeAt image j i} for all 84 words, each
  by \lean{rfl}, so the kernel decodes every concrete word; the file takes 3.4\,s. A mismatch between
  the disassembly and the contract decoder would fail the build.
\item Block proofs are forward chains \lean{t.trans (Reach.xxx code\_j rfl …)}. States are
  abstracted at block boundaries.
\item Loops are proved by list induction: \lean{loop\_run} over a segment's pairs, \lean{rows\_run} over
  the rows.
\item The 19-instruction segment is proved once, as \lean{SegCode j res}, and instantiated at
  $j=14,33,52$.
\end{itemize}
No symbolic executor was needed at this image size. In v0.2 the machine proofs took about 3.3 lines per
instruction, including loops and composition; A3's calibration was 8.5 lines per instruction for a
block alone.

\begin{table}[t]
\centering\footnotesize
\setlength{\tabcolsep}{4pt}
\begin{tabular}{@{}p{0.36\textwidth}lrrrr@{}}
\toprule
Component & Files & v0.2 & v0.3 & $\Delta$ & A3 estimate \\
\midrule
Spec: layout, sparse encoding, walk, CSR lemma, protocol, security & \code{Encoding}, \code{Protocol} & 319 & 328 & $+9$ & 400--750 \\
RISC-V step-lemma tooling (reusable; a \textsf{lib-L0} prototype): counters, footprint, potential & \code{Rv/L0} & 171 & 339 & $+168$ & (library) \\
Field arithmetic bridge & \code{Arith} & 119 & 119 & 0 & 50--100 \\
Machine proofs: segment, row, prologue and run, closed forms, first touch, latency & \code{Seg}, \code{Row}, \code{Run}, \code{Main} & 705 & 886 & $+181$ & 1{,}000--1{,}850 \\
Generated: image literal, decode table & \code{Image}, \code{Code} & 108 & 101 & $-7$ & --- \\
Certificate, score, solution(s), root & \code{Certificate}, \code{Score}, \dots & 59 & 147 & $+88$ & 50--100 \\
\midrule
\textbf{Total} (theorems: 244 / 311) & & \textbf{1{,}481} & \textbf{1{,}920} & $\mathbf{+439}$ & 1.5--2.8k + library \\
\bottomrule
\end{tabular}
\caption{Proof size in code lines, counted by \code{loc.py} [P3~§4.3; P3b~§4; \repo{logs/loc.log}].
v0.3 by file: \code{Encoding} 159, \code{Protocol} 169, \code{Seg} 274, \code{Row} 155, \code{Run}
239, \code{Main} 218, \code{Certificate} 56, \code{Score} 34, \code{BaselineSolution} 25,
\code{BaselineSolutionL1} 30, root 2. The A3 column sums A3~§3.1's rows ``with \textsf{lib-L0}--\textsf{L2}'': rows 1--2 for
the spec, rows 3--5 and 7--8 for the machine. A3's estimate without a library was
2.7--4.5k.}\label{tab:base-loc}
\end{table}

\Cref{tab:base-loc} gives the line counts. The migration added 439 lines. About 170 of them are the
latency clause (the potential in \code{Rv/L0.lean}, \lean{rowsLat\_eq}, \lean{baseline\_latency}, the latency record and
its 30-line solution file). Without latency the increase would be about 270, inside the v0.3 design's
pre-registered estimate of $+150$ to $+300$ [P3b~§0]. The total remains at the low end of A3's
with-library range, although no organiser library existed and P3 and P3b wrote their own \textsf{lib-L0} prototype.

\paragraph{A finding: nanoda runs out of memory on \lean{simp} rewriting under large machine
states.} The first harness run with nanoda as the second kernel was killed by the operating system
at 60\,GB RSS (\code{journalctl}: \code{Killed process … (nanoda\_bin) anon-rss:60768196kB}). P3b
bisected it with capped probes, exporting the closure of single constants under \code{ulimit -v}.
\lean{seg\_run}, \lean{prologue} and \lean{baseline\_secure} (64\,s, 2.3\,GB) checked fine, but
\lean{row\_run} diverged. Inside it the culprit was the accept branch, and inside that a single goal:
\lean{s'.mem = M}, proved by \lean{simp [m3]}. Simp's rewrite proof runs through the nested
\lean{setReg}/\lean{setPC} state, which contains the branch \lean{if} over \lean{clResult} terms, and
nanoda's check of it explodes: more than 6\,GB in 100\,s. \lean{exact m3} is definitional and checks
in 1.7\,s. After that one-line change, nanoda checks the whole certificate closure (a 127\,MB export)
in 62\,s and 2.4\,GB [P3b~§5, §7]. The Lean default kernel had accepted the \lean{simp} proof
throughout. \emph{Lesson:} a second kernel is a resource risk as well as a trust gain, and the risk
sits in proof style, not in the statement. The organiser library should ship definitional frame
lemmas, for example a \lean{mem\_after\_steps}, and a lint that warns on \lean{simp} rewriting under
large machine states, so that submitters do not meet a 60\,GB OOM in CI [P3b~§8 item 1].

\paragraph{Proof traps, as lessons for the organiser library.} \Cref{tab:base-traps} lists the traps
A3 predicted, the four new ones P3 met in v0.2 and the nanoda trap of v0.3. Each has a library-level
remedy.

\begingroup\footnotesize
\setlength{\tabcolsep}{4pt}
\begin{longtable}{@{}>{\raggedright\arraybackslash}p{0.22\textwidth}>{\raggedright\arraybackslash}p{0.06\textwidth}>{\raggedright\arraybackslash}p{0.30\textwidth}>{\raggedright\arraybackslash}p{0.34\textwidth}@{}}
\caption{Proof traps [A3-C3; P3~§7; P3b~§5]. ``A3'' marks a trap predicted before the baseline was
written, ``new'' one found while writing v0.2, ``v0.3'' one found during the migration.}\label{tab:base-traps}\\
\toprule
Trap & Origin & Symptom & Remedy (library action) \\
\midrule
\endfirsthead
\toprule
Trap & Origin & Symptom & Remedy (library action) \\
\midrule
\endhead
Naive unfolding of \lean{execute} & A3 & \lean{simp only [run, execInstrBr, …]} over 12 steps does not finish in 5\,min & never unfold a long run; per-instruction \lean{Reach} lemmas composed by \lean{Reach.trans} (\textsf{lib-L0}) \\
Concrete \code{jal} offsets & A3 & \lean{rfl} hits max recursion; kernel ``deep recursion'' after 266\,s & avoided: no \code{jal}; branch offsets enter only through \lean{signExtend13 o = …} facts proved once by \lean{rfl}; \textsf{lib-L0} must state \code{jal}/\code{jalr} with symbolic offsets \\
\lean{decide} on \lean{BitVec} constants & A3 & kernel deep recursion & \lean{rfl} for sign-extension constants \\
\texttt{QUERY} memory in nested states & new & the closed-form memory lambda is copied into every later state term: error messages of thousands of lines, slower elaboration & abstract the state (\lean{obtain ⟨s', …⟩}) right after each \texttt{QUERY}; ship readback lemmas over an opaque post-state \\
\lean{omega} on nested \lean{\% 2\^{}64} & new & hangs for more than 60\,s, e.g.\ on \lean{ofNat n + ofNat (2\^{}64−1) = ofNat (n−1)} & explicit \lean{Nat.mod\_eq\_of\_lt}/\lean{Nat.add\_mod\_right} rewriting; ship word-arithmetic normal forms \\
\lean{simp} on layout constants & new & kernel deep recursion from \lean{simp [XB, WB, TAPE\_BASE, csz, classSizes]} proving $2^{32}+32+8\cdot2^{10}\le2^{33}$ & a standalone \lean{norm\_num} lemma on the unfolded constants; ship the memory-map inequalities as lemmas \\
\lean{rfl} on the honest transcript & new & \lean{prover.transcript … 0 = msgOf w} hits the elaborator's max recursion depth at class sizes & unfold \lean{Strategy.prefixes} and \lean{Fin.induction} with \lean{simp only}, then \lean{simp}; ship a transcript-of-honest-prover lemma \\
\lean{simp} rewriting under a large machine state & v0.3 & accepted by the Lean kernel, but nanoda OOM at 60\,GB (one goal: more than 6\,GB in 100\,s) & a definitional \lean{exact} (1.7\,s); ship frame lemmas such as \lean{mem\_after\_steps}; run nanoda in CI with a memory cap \\
\bottomrule
\end{longtable}
\endgroup

Two smaller frictions recur. \lean{decide\_eq\_true\_eq} does not fire on \lean{decide (a ≤ b) =
true} produced by the machine's \lean{writable}/\lean{inMemory} checks; \lean{decide\_eq\_true} has
to be applied explicitly. And \lean{omega} fails on goals containing the \lean{True ∧ …} conjuncts
that \lean{mk.injEq} leaves behind; \lean{and\_intros <;> first | trivial | omega} works.

Recommendations for the organiser library (\cref{sec:issues}) [P3~§8; P3b~§8]: ship \code{Rv/L0.lean}
as the \textsf{lib-L0} starter, extended to stores, W-ops, \code{jal}/\code{jalr} and \code{mul}; ship the
code-table generator and image renderer (\repo{asm/gen\_code.py}, \repo{asm/gen\_image.py}) and
\code{rfemu} with its differential test against \lean{runFull}; add a counted-loop combinator (about 90
lines of \lean{loop\_run}/\lean{rows\_run} boilerplate), definitional frame lemmas and a nanoda-in-CI
step; move \code{Arith.lean} into \textsf{lib-L2}; and keep the symbolic executor on the plan for
succinct entries, which will have thousands of instructions rather than 84.

%-----------------------------------------------------------------------------------------------
\subsection{Check time and resources}\label{subsec:base-time}

\Cref{tab:base-time} gives the check measurements. The machine had 32 cores and 62\,GB, shared with
other agents, so times vary by about $\pm15\%$. The toolchain was Lean and Mathlib v4.34.1
\cite{lean4340}. The v0.3 harness runs the comparator \cite{comparator,lean4export} and the gate under
the contract's bubblewrap landrun shim (\code{bwrap-landrun.sh}); production must use the real
sandbox [P3b~§9].

\begin{table}[t]
\centering\footnotesize
\begin{tabular}{@{}>{\raggedright\arraybackslash}p{0.47\textwidth}>{\raggedright\arraybackslash}p{0.24\textwidth}>{\raggedright\arraybackslash}p{0.17\textwidth}@{}}
\toprule
Measurement & Wall & Peak RSS \\
\midrule
\multicolumn{3}{@{}l}{\emph{v0.3 (current) [P3b~§7]}} \\
Full warm \code{scripts/check.sh}: image, claim regeneration from Lean, hygiene, render $\times2$, harness $\times2$ with both kernels, build, axioms, \code{leanchecker} $\times2$, differential test & 814.8\,s [P3b] & 3.5\,GB \\
Harness T1, stages 0--2 (two fresh-tree builds of the candidate), Lean kernel + nanoda & 369.0\,s; comparator stage 243.6\,s & 3.5\,GB \\
Harness \lean{Latency1} record, the same & 363.0\,s; comparator stage 238.0\,s & 3.5\,GB \\
Harness T1 / \lean{Latency1}, default kernel only (first runs) & 340.6\,s / 310.9\,s & 3.5\,GB \\
Organiser package from scratch (\code{org-build.sh}; Mathlib oleans present) & 110.5\,s (247\,s user) & 3.5\,GB \\
nanoda alone on the certificate closure (127\,MB export) & 62\,s & 2.4\,GB \\
Gate / \code{leanchecker} (each solution) & 10.7\,s / 24\,s & 2.3\,GB \\
Independent re-run of \code{check.sh} [O-recheck-v03]: \code{CHECK OK}, exit 0; gate closures 24{,}388 / 24{,}406 constants & 13\,min 35\,s (7\,min 56\,s user) & 3.5\,GB \\
\quad of which harness T1 / \lean{Latency1}, both kernels: \code{HARNESS: ACCEPT} & 365.8\,s / 373.2\,s & 3.52 / 3.51\,GB \\
\midrule
\multicolumn{3}{@{}l}{\emph{v0.2 (historical) [P3~§0, §4.2; O-recheck]}} \\
Full warm check (no harness; default kernel only) & 109\,s [P3]; 135\,s [O-recheck] & 3.5\,GB \\
Project-cold: cold build + comparator & $\approx3.7$\,min & 3.5\,GB \\
With a fresh Mathlib download (P1: 7.8\,min clean setup) & $\approx12$\,min & 3.5\,GB \\
\midrule
Brief's requirement & $<1$\,h & \\
\bottomrule
\end{tabular}
\caption{Check time and resources [P3b~§5, §7; \repo{logs/check.log}, \repo{harness-T1.log},
\repo{harness-L1.log}, \repo{gate.log}, \repo{leanchecker.log}; P3~§4.2; O-recheck]. A3's
pre-registered estimates were about 8\,min warm and 20--25\,min cold [A3~§3.2].}\label{tab:base-time}
\end{table}

\paragraph{Against the brief's one hour.} The organiser's cost to accept one record is one harness
run, about 365\,s, or about 8\,min including a from-scratch organiser package (111\,s). The package is
built once and shared read-only by every harness run. The submitter's full warm check, which runs the harness
for both records with both kernels, takes 815\,s, about 14\,min. Both are well inside the brief's
1\,h, by factors of about 10 and 4. The independent re-run took 13\,min 35\,s [O-recheck-v03]. The v0.3 check takes 6--7.5 times as long as v0.2's warm check. The harness
order explains most of it: two fresh-tree builds of the candidate per run, a second kernel, and two
records. Per-module elaboration ranges from 2.4\,s (\code{Image}) to about 10\,s (\code{Row})
[P3b~§7]; P3 measured an import floor of about 2.5\,s per module. The emulator measurements are not part of
the check: at $k=24$ a run needs about 10\,GB of RAM and about 40\,min including instance generation.

\paragraph{Reproduce.} From the repository root:
\begin{Verbatim}[fontsize=\small,frame=single,framerule=0.3pt,rulecolor=\color{black!30},framesep=4pt]
contract-v03/scripts/org-build.sh contract-v03/.lake/org  # organiser package
submissions/baseline-v03/scripts/check.sh      # all checks, harness T1 + Latency1
submissions/baseline-v03/emu/measure.sh 16 20 24   # class runs (k=24: ~10 GB, ~40 min)
\end{Verbatim}
\code{check.sh} prints \code{HARNESS: ACCEPT} per record, the differential-test \code{MATCH} lines and
\code{CHECK OK}. It needs the assembler \code{riscv64-unknown-elf-as}. The only edit to
\repo{contract-v03/} is in \repo{lakefile.toml}: the existing \lean{RecursionFastCandidate} lib now
has its source directory in \repo{submissions/baseline-v03}, and a new lib
\lean{RecursionFastBaselineAudit} holds the rendered challenges and solutions. Neither is a default
target, and \code{org-build.sh} drops both [P3b~§10]. The v0.2 submission is reproduced with
\code{bash submissions/baseline/scripts/check.sh} against \repo{contract/} (\cref{sec:reproducibility}).

%-----------------------------------------------------------------------------------------------
\subsection{The historical v0.2 certification}\label{subsec:base-v02}

The first certificate, P3's, is against the trusted contract v0.2; it is kept intact in the directory
\repo{submissions/baseline/} [P3]. Its theorem \lean{certificate}, of type
\lean{Certificate1 submission claim}, covers
security, admission, image size, \lean{Agrees} and \lean{AcceptingCostAtMost} over v0.2's ten
counters ($\mathrm{alu}=13+18m+3P$, $\mathrm{clmul}=14m+2P$, $\mathrm{load}=1+3m+3P$,
$\mathrm{ctrl}=5m+P$, $\mathrm{sys}=1$); comparator replayed it in the default kernel, and the
independent re-check reproduced it [O-recheck]. Under the flat v0.2 weights it scored
$S_k=452\cdot2^k+1{,}093$ exactly, i.e.\ $\mathrm{Score}_{\mathrm{T1}}=4.74\times10^8$ [P3~§5], 6.3\%
below D3's pre-registered model of $5.06\times10^8$ [D3~§5.3]: storing $z$-addresses makes the
instructions cheaper than modelled ($322\cdot2^k$ against $454\cdot2^k$), while whole power-of-two
leaves make the bytes dearer ($130\cdot2^k$ against $\approx28.5\cdot2^k$; I-BAS-2). The baseline
was tuned neither to the v0.2 weights nor to the v0.3 profiles, as the brief asks. The ``trivial
baseline (modelled)'' rows of \cref{tab:bench-trajectory,tab:bench-models} are D3's modelled
baseline and D4's re-pricing of it; the certified image is priced in \cref{tab:base-score}.

%-----------------------------------------------------------------------------------------------
\subsection{Contract findings and the v0.2\texorpdfstring{$\to$}{ to }v0.3 migration}\label{subsec:base-findings}

Writing the baseline was the first end-to-end use of the contract. The v0.2 submission produced seven
findings [P3~§6], and the v0.3 migration five more [P3b~§8]. I-BAS-2, the actionable one still open,
is in \cref{subsec:base-score}.

\begin{issue}{I-BAS-1}{The dense \lean{Index} made the honest indexer infeasible to evaluate (resolved in v0.3)}
In v0.2, \lean{Index sz} stored \lean{Mat sz = Fin m → Fin cols → K}. Any computable \lean{IDX} had
to scan all $m\cdot\mathrm{cols}=2^{2k}$ entries to find the support, which is $2^{48}$ at $k=24$, so
the badge path (\cref{subsec:bench-badge}) could not run the compiled indexer on the benchmark
indices. \emph{Status:} resolved. v0.3's \lean{Index} is sparse (entry lists with
\lean{nnz\_le}), and the baseline's indexer is the one-pass \lean{buckets} with \lean{getD\_buckets}
(\cref{subsec:base-spec}); see row V37 of \cref{tab:iss-resolved}. It has not yet been run at class
size.
\end{issue}

The remaining findings, with their v0.3 status:
\begin{itemize}[nosep]
\item \emph{Ignored witness entries.} $n_w=2^k-1025$ is not a power of two, so the message carries
  1{,}025 entries the verifier ignores. This is harmless, since \lean{Agrees} quantifies over every raw
  message, ignored entries included. The relaxation proposed in I-BAS-2 removes it.
\item \emph{\lean{Instruction} had no \lean{DecidableEq}.} v0.3 derives it; \lean{rfl} remains the
  fast way to prove the \lean{CodeAt} table.
\item \emph{Rendered image literals.} v0.3 renders the image (F2). The baseline now emits it in the
  renderer's chunked form, so \lean{Challenge.image = image} is \lean{rfl}, and the \lean{CodeAt} table
  is unchanged. The renderer needs the image files next to the claim
  (\repo{audit/baseline.bin}, \repo{audit/baseline.data} are copies) [P3b~§8 item 4].
\item \emph{Counts versus timing.} Linear re-weighting stays free, since the per-block counts are
  exact. The latency clause needed new proof, as predicted: the timer potential in \code{Rv/L0.lean}, about 170
  lines in all, with 1.46 times slack against the scoreboard.
\item \emph{The tape header carries the class cap $\mathsf{nnz}$, not the actual count.} This is
  unchanged in v0.3 and harmless: the baseline reads end words from the index.
\item \emph{The \lean{V02} transfer was not needed} [P3b~§8 item 3]. A single reachability predicate
  on the contract's \lean{execute}, with an accounting invariant and a timer potential, gives
  \lean{Agrees}, cost and latency at once. \lean{Seq.run\_eq} is the cheaper path if the latency
  board is skipped. The library's \code{V02.lean} is therefore insurance rather than the expected
  migration route.
\item \emph{A second kernel can fail on resources} (the nanoda finding of \cref{subsec:base-proof}).
\end{itemize}

\paragraph{What the migration did.} Measured against the v0.3 migration plan (DESIGN §15), the
baseline adopted the sparse \lean{Index}, all 18 counters with first touch, the latency clause, the
rendered image, claim and challenge format (F2, F4, F5) with claimed data generated from the
contract's Lean functions, the absorption and leaf-hash pricing (F6, P1), and the round-3 harness and
hash-manifest gate (G0, G1/G2, F3). The image did not change. The migration cost 439 lines, and the
only proof-style change forced by the harness was the nanoda fix.

%-----------------------------------------------------------------------------------------------
\subsection{How agents are expected to beat it}\label{subsec:base-beat}

\paragraph{Within the trivial family: constant factors only.} The claimed bound must hold for
\emph{every} index of the class shape (\cref{subsec:bench-role}), so benchmark structure cannot lower
it. Exploiting coefficient-1 entries or the actual nonzero count does not help, because the
worst-case index has neither. What remains:
\begin{itemize}[nosep]
\item removing the index padding of I-BAS-2, about 39\% of the RISC~Zero score and less in the
  profiles that do not price pages;
\item shaving per-row and per-nonzero instructions;
\item re-laying out the index;
\item for the latency board, a tighter timing proof (up to 1.46 times) or harts.
\end{itemize}
None of these changes the $\Theta(\mathsf{nnz})$ instruction, byte and page cost of reading the whole
index.

\paragraph{Succinct designs: two orders of magnitude.} The expected trajectory is in
\cref{subsec:bench-trajectory} (\cref{tab:bench-trajectory}). The first proven-regime succinct design
there, B1, is Spartan with a binary WHIR/Ligerito opening. Under the v0.2 weights it is modelled at
$\mathrm{Score}\approx5.26$\,M. That is about 90 times below the baseline's v0.2-weight geomean of
474\,M (\cref{tab:base-score}, last column) and 96 times below D3's modelled baseline. Those rows are
model estimates accurate to about a factor of 2, not certified numbers, and they have not been
re-priced under the final round-3 pricing. D4's re-run under the pre-round-3 profiles put B1 at
4.7--10.3\,M per profile (\cref{tab:bench-models}), against the certified baseline's per-profile
geomeans of 844\,M--2.36\,G (\cref{tab:base-score}). The ratios, about 180--280, are indicative
only, because the two sides use different leaf pricing. The rows point to holography engineering as the lever: query count, leaf size and
the size of the index commitment.

\paragraph{The hybrid small-class fallback.} A \lean{Submission1} has \emph{one} image but a protocol
per class (\lean{protocol : (c : Fin 3) → Protocol (classSz c)}). An image can therefore dispatch on
the tape header $(m,n_x,n_w,\mathsf{nnz})$ and run the baseline's check at the classes where a
succinct design costs more. D2 expected this hybrid to be the first improvement under prime fields,
where the crossover was $k\approx16$. Over binary fields D3's model puts the crossover at $k\approx13$
[D3~§6.2], so a mature B1-class design has no use for the fallback at $2^{16}$. An early succinct
entry, for example one in the unique-decoding regime or with an untuned opening, may still exceed the
baseline's $S_{P,16}$ of $5.3\times10^7$--$1.5\times10^8$ (\cref{tab:base-score}). The fallback caps
its $2^{16}$ score at the baseline's.

Because the score is a geometric mean, the gain is at most a cube root. Taking the fallback means
proving the baseline branch inside the hybrid image. \code{Rv/L0.lean} and the spec are reusable as
they stand. The block lemmas are tied to \lean{image} through the generated \lean{CodeAt} table and
\lean{SegCode}, so they must be re-targeted at the new code positions. That is mechanical work over
the 886 machine-proof lines, not new mathematics.
\endgroup
